\subsection{Exactly when an accessible neighbourhood kills a new residue}

\begin{claim}{D.28}{Proved}\label{clm:D28}
Let $k$ be a number field, $U/k$ a smooth geometrically connected
hyperbolic curve, $s:G_k\to E_U$ a continuous section, and $a\in U(k)$
an ordinary point. Choose actual sections and put
$S=s(G_k)$, $D_a=s_a(G_k)$ and $\Pi=\pi_1(U_{\bar k})$.
For a continuous character $\chi:G_k\to\mathbf Z/n$, define the
compact subset
\[
 C_\chi=\{\gamma s_a(\sigma)\gamma^{-1}:
       \gamma\in\Pi,\ \chi(\sigma)\ne0\}\subset E_U.
                                                               \tag{D.28.1}
\]
There is a finite \'etale neighbourhood $V\to U$ carrying $s$ such
that $\chi$ becomes trivial over every residue field above $a$ if
and only if
\[
                         S\cap C_\chi=\varnothing.             \tag{D.28.2}
\]
Equivalently, no $\sigma$ with $\chi(\sigma)\ne0$ has
$s(\sigma)$ geometrically conjugate to $s_a(\sigma)$.
This is an elementwise condition; nonconjugacy of the two sections
as homomorphisms is insufficient.

For $U=\mathbf P^1_k-\{0,1,\infty\}$, $a\in k-\{0,1\}$,
condition (D.28.2) supplies the canonical reciprocity test for
$\kappa(t-a)\smile\chi$ in D.27, using the section on $V$.
The assertion that such a neighbourhood always exists is false,
even for an actual rational point section: take
\[
       k=\Q,\qquad s=s_3,\qquad a=2,\qquad
       \chi=\chi_{\Q(i)/\Q}.
\]
The canonical value in this example is $3-2=1$, so the reciprocity
test itself nevertheless vanishes.
\end{claim}
\begin{proof}
A finite \'etale neighbourhood of $s$ corresponds to an open subgroup
$H\subset E_U$ containing $S$. It is geometrically connected:
$H\to G_k$ is surjective, and consequently the geometric group acts
transitively on the finite coset set. The fibre over $a$ is the finite
$D_a$-set of cosets $E_U/H$. Every coset has a representative in $\Pi$,
since $E_U=\Pi H$. Its stabilizers are
\[
                       D_a\cap\gamma H\gamma^{-1}.
\]
Under $D_a\simeq G_k$, these are the absolute Galois groups of the
residue fields of the points in the fibre, with their selected
embeddings. This is the finite \'etale Galois-category correspondence
\citep[Expos\'e~V]{E_SGA1}. Thus $\chi$ vanishes on all these residue
subgroups exactly when $H\cap C_\chi$ is empty. Necessity of
(D.28.2) follows from $S\subset H$.

Conversely, $C_\chi$ is compact: it is the continuous image of the
product of $\Pi$ and the clopen set $G_k-\ker\chi$. Let $N$ run through
open normal subgroups of $E_U$. Each $SN$ is an open subgroup containing
$S$, and
\[
                              \bigcap_N SN=S,
\]
since $S$ is closed. If every $SN$ met $C_\chi$, the closed subsets
$SN\cap C_\chi$ would have the finite-intersection property: the
intersection of finitely many $N$ refines all the corresponding
conditions. Compactness would then give a point of $S\cap C_\chi$.
Under (D.28.2) this is impossible, so some $H=SN$ is the required
neighbourhood. This uses compactness of one fixed conjugacy set,
not compactness of a space of continuous sections, and does not
require finite generation.

If $s(\tau)=\gamma s_a(\sigma)\gamma^{-1}$, projecting to $G_k$
gives $\tau=\sigma$, because $\gamma\in\Pi$. This proves the
stated equivalent elementwise formulation.

\emph{The positive reciprocity consequence.}
Now use the three-punctured line and put $W=U-\{a\}$. Let
$V_a\subset V$ be the finite fibre above $a$ and write
\[
          \beta=\kappa(t-a)\smile\chi\in\operatorname{Br}(W)[n].
\]
The residue of its pullback to $V-V_a$ at $y\in V_a$ is
\[
                         \operatorname{Res}_{k(y)/k}\chi.
\]
The ramification index is one, because $V\to U$ is \'etale.
All these residues vanish by construction. Purity for a regular curve
therefore extends the pulled-back class uniquely to
$\operatorname{Br}(V)[n]$. We use the cohomological residue convention
of D.27, the negative of the Witt convention; the zero-residue
criterion is unchanged
\citep[\S1.4.1 and Proposition~3.6.1]{D_ColliotTheleneSkorobogatov2019}.

Choose a lift of that Brauer class to $H^2(V,\mu_n)$. The affine
hyperbolic curve is an \'etale $K(\pi,1)$, so the lifted section $s_V$
pulls it back to $H^2(k,\mu_n)=\operatorname{Br}(k)[n]$.
The ambiguity lies in $\operatorname{Pic}(V)/n$; D.23's Chern-class
vanishing kills it. Equivalently, the local point/cusp pullbacks
vanish on that ambiguity, and Brauer--Hasse--Noether kills it globally.
Global reciprocity applies to the resulting class.

At an ordinary local point $x_v\ne a$, the local lifted section is
an ordinary point of $V(k_v)$ above $x_v$, by D.6 and D.19 and
functoriality of the point sections. Its evaluation equals
$\beta(x_v)$. If $x_v=a$, its rational lift to the finite fibre gives
an embedding $k(y)\hookrightarrow k_v$ for some $y\in V_a$.
Consequently $\chi_v=0$; any separately chosen new tangential class
has zero pairing with $\chi_v$. At an old cusp, functoriality gives
the already specified old tangential leading-term evaluation, as in
D.27. The sum of these local invariants is therefore zero. It is a
finite sum because it is the localization of one global Brauer class.

This proves one actual missing canonical test whenever (D.28.2)
holds. It does not establish restricted-product admissibility for
the translated Kummer tuple. The finite support just obtained is
support of the pairing with this fixed $\chi$, and can depend on
$\chi$ and the chosen neighbourhood. Nor does it prove (D.28.2)
for every character.

\emph{The real obstruction to universal residue killing.}
Take the data in the final part of the statement. Every finite
\'etale neighbourhood $V\to U$ carrying $s_3$ has a $\Q$-rational
point above $3$: the section fixes a geometric point of that finite
fibre. On real smooth loci a finite \'etale map is a finite
topological covering. Indeed it is locally a homeomorphism by the
implicit function theorem and is proper, since the algebraic map
is finite. Lift the path from $3$ to $2$ inside $(1,\infty)$,
starting at that real point. Its endpoint is a real point above $2$.

The corresponding closed residue field $k(y)$ therefore has a real
embedding. It cannot contain $\Q(i)$, so the restriction of $\chi$
to $G_{k(y)}$ is nonzero. Thus no accessible neighbourhood kills
$\chi$ on every point of the fibre. In the equivalent path formulation,
the real path gives a geometric \'etale path fixed by complex
conjugation $c$, hence
\[
       s_3(c)=\gamma s_2(c)\gamma^{-1},\qquad\chi(c)\ne0.
\]
It produces an element of $S\cap C_\chi$.

There is no failure of reciprocity here: the actual global value is
$1$, whose Kummer class and every cup product vanish. The example
only disproves the proposed universal method of obtaining that test
by killing every residue on one accessible finite cover. The same
obstruction occurs over any real place when the two points lie in
the same real component and the character is nontrivial on its
complex conjugation. Restricting to even characters or passing to
a totally imaginary field removes this example, without proving
(D.28.2).
\end{proof}

C.7h does not close this gap. Its primitive-valuation uniqueness
separates distinct ordinary points when their sections are compared
over a suitable whole local subgroup with nonzero tame
$\ell$-image. An individual global element occurring in (D.28.1)
need not generate such a subgroup. The difference between one
conjugator for whole sections and separate conjugators for individual
elements remains essential. D.28 neither excludes another
reciprocity argument nor proves the new-puncture restricted-product
condition, and leaves D1--D4 unresolved.
